\section{Evaluation}
\label{sec:evaluation}

\subsection{Experimental Setup}

We apply AreaSection to all 50 states using 2020 Census data, TIGER
tract geometries, and 2020 presidential election results (precinct-to-tract
interpolation via VEST/Fekrazad \citep{fekrazad2021}).
All runs use 50 seeds per ratio.
Population balance tolerance is 1.5\,\% (broader than the 0.5\,\%
constitutional standard, applied to accommodate discrete-tract rounding in
recursive bisections; see \S\ref{sec:limitations}).

The reference implementation is the \texttt{redist} Rust binary using METIS 5.2.1
compiled for Windows x64, linked via the \texttt{metis-rs} FFI crate.
44 of 50 states converge within tolerance; 6 states (FL, IL, MI, NY, PA, TX)
exceed the 1.5\,\% balance threshold due to recursive bisection rounding in
subgraphs created by asymmetric first-level splits (see \S\ref{sec:limitations}).

\subsection{Natural Ratio Results: 50-State Sweep}

Table~\ref{tab:natural-ratios} reports the natural ratio selected by AreaSection
for each state, compared to the GeoSection natural ratio.
The contrast is stark.

\begin{table}[h]
\centering
\caption{AreaSection vs GeoSection natural ratios, 50 states.
  ``Equal/near-equal'' = ratio within 10\,\% of $k/2:k/2$.
  EC = edge cut at natural ratio.}
\label{tab:natural-ratios}
\small
\begin{tabular}{llrrrrl}
\toprule
State & Seats & GeoSection & AS ratio & AS EC (km) & $p^*$ & Category \\
\midrule
FL  & 28 & 1:27 & 14:14 & 369 & -- & Equal \\
NY  & 26 & 1:25 & 13:13 & 402 & -- & Equal \\
IL  & 17 & 1:16 & 8:9   & 568 & -- & Near-equal \\
OH  & 15 & 1:14 & 7:8   & 456 & -- & Near-equal \\
PA  & 17 & 1:16 & 7:10  & 375 & -- & Near-equal \\
MI  & 13 & 1:12 & 5:8   & 580 & -- & Near-equal \\
TX  & 38 & 1:37 & 13:25 & 930 & -- & Asymmetric \\
CA  & 52 & 1:51 & 20:32 & 533 & -- & Asymmetric \\
NC  & 14 & 1:13 & 6:8   & 284 & -- & Near-equal \\
GA  & 14 & 1:13 & 3:11  & 483 & 0.91 & Concentrated \\
WI  &  8 & 1:7  & 3:5   & 622 & 0.93 & Asymmetric \\  % sweep used 4:4 but 50-seed result
VA  & 11 & 1:10 & 5:6   & 437 & -- & Near-equal \\
AZ  &  9 & 1:8  & 4:5   & 851 & -- & Near-equal \\
IN  &  9 & 1:8  & 4:5   & 306 & -- & Near-equal \\
TN  &  9 & 1:8  & 4:5   & 274 & -- & Near-equal \\
NJ  & 12 & 1:11 & 4:8   & 97  & 0.82 & Concentrated \\
MA  &  9 & 1:8  & 2:7   & 95  & 0.89 & Concentrated \\
MN  &  8 & 1:7  & 4:4   & 573 & -- & Equal \\
MO  &  8 & 1:7  & 4:4   & 598 & -- & Equal \\
CO  &  8 & 1:7  & 4:4   & 664 & -- & Equal \\
MD  &  8 & 1:7  & 4:4   & 179 & -- & Equal \\
WA  & 10 & 1:9  & 5:5   & 732 & -- & Equal \\
KY  &  6 & 1:5  & 3:3   & 317 & -- & Equal \\
OR  &  6 & 1:5  & 3:3   & 823 & -- & Equal \\
AL  &  7 & 1:6  & 3:4   & 534 & -- & Near-equal \\
SC  &  7 & 1:6  & 3:4   & 334 & -- & Near-equal \\
LA  &  6 & 1:5  & 2:4   & 350 & -- & Asymmetric \\
OK  &  5 & 1:4  & 2:3   & 438 & -- & Near-equal \\
CT  &  5 & 1:4  & 2:3   & 139 & -- & Near-equal \\
KS  &  4 & 1:3  & 1:3   & 418 & 0.78 & Concentrated \\
MS  &  4 & 1:3  & 2:2   & 404 & -- & Equal \\
IA  &  4 & 1:3  & 2:2   & 414 & -- & Equal \\
AR  &  4 & 1:3  & 2:2   & 552 & -- & Equal \\
UT  &  4 & 1:3  & 2:2   & 753 & -- & Equal \\
NV  &  4 & 1:3  & 2:2   & 731 & -- & Equal \\
NM  &  3 & 1:2  & 1:2   & 706 & -- & Unchanged \\
NE  &  3 & 1:2  & 1:2   & 504 & -- & Unchanged \\
HI  &  2 & 1:1  & 1:1   & 79  & -- & Equal \\
ID  &  2 & 1:1  & 1:1   & 381 & -- & Equal \\
ME  &  2 & 1:1  & 1:1   & 228 & -- & Equal \\
MT  &  2 & 1:1  & 1:1   & 753 & -- & Equal \\
NH  &  2 & 1:1  & 1:1   & 128 & -- & Equal \\
RI  &  2 & 1:1  & 1:1   & 51  & -- & Equal \\
WV  &  2 & 1:1  & 1:1   & 349 & -- & Equal \\
AK  &  1 & --   & --    & --  & -- & Single \\
DE  &  1 & --   & --    & --  & -- & Single \\
ND  &  1 & --   & --    & --  & -- & Single \\
SD  &  1 & --   & --    & --  & -- & Single \\
VT  &  1 & --   & --    & --  & -- & Single \\
WY  &  1 & --   & --    & --  & -- & Single \\
\bottomrule
\end{tabular}
\end{table}

\noindent\textbf{Pattern 1: Large equal-population states force equal splits.}
Florida (28 seats), New York (26 seats), and Minnesota/Missouri/Colorado/Maryland
(8 seats each) all receive exactly equal 50/50 splits.
This is not a coincidence: for states with relatively uniform density, the
50\,\% area constraint is naturally compatible with a 50/50 population split.

\noindent\textbf{Pattern 2: Urban-dominant states resist equal splits but
improve dramatically.}
Illinois moves from 1:16 (GeoSection) to 8:9 (AreaSection).
Texas moves from 1:37 to 13:25.
California moves from 1:51 to 20:32.
The area constraint cannot force a perfectly equal split because Chicago,
Houston, and Los Angeles each occupy a small fraction of their state's land area,
making the dual constraint infeasible at equal population ratios.
But the improvement over GeoSection is substantial.

\noindent\textbf{Pattern 3: Highly concentrated metros retain asymmetric splits.}
Georgia (3:11), Massachusetts (2:7), and New Jersey (4:8) retain strongly
asymmetric splits.
The Lorenz analysis confirms why: Georgia's dense half (sorted by density)
holds $p^* \approx 0.91$ of the state's population.
Achieving a 7:7 population split while maintaining 50\,\% land area is
geometrically infeasible: any contiguous 50\,\% of Georgia's land either
captures $<30$\,\% of the population (the rural half) or $>70$\,\% (the
Atlanta half), with no geography in between.

\subsection{Wisconsin: Head-to-Head Partisan Comparison}

Table~\ref{tab:wi-partisan} compares GeoSection and AreaSection on Wisconsin
(8 districts, 2020 presidential results).

\begin{table}[h]
\centering
\caption{District-level Democratic vote share: GeoSection vs AreaSection,
Wisconsin 2020. Sorted by district index.}
\label{tab:wi-partisan}
\begin{tabular}{rrrr}
\toprule
District & GeoSection D\,\% & AreaSection D\,\% & Category \\
\midrule
1 & 73.6 & 72.8 & Safe D \\
2 & 49.4 & 43.0 & Near-toss-up / Safe R \\
3 & 49.4 & 51.5 & Near-toss-up / Lean D \\
4 & 37.8 & 42.7 & Safe R \\
5 & 53.4 & 68.5 & Lean D / Safe D \\
6 & 58.7 & 42.1 & Lean D / Safe R \\
7 & 44.6 & 38.4 & Safe R \\
8 & 40.3 & 46.2 & Safe R / Lean R \\
\midrule
D seats & 3 & 3 & Same \\
R seats & 5 & 5 & Same \\
D vote share & 50.3\,\% & 50.3\,\% & Same \\
Prop.\ gap & $-12.8$\,pp & $-12.8$\,pp & Same \\
\bottomrule
\end{tabular}
\end{table}

\noindent\textbf{Key finding: identical seats, different competitiveness.}
Both algorithms produce 3 Democratic seats and 5 Republican seats, with a
$-12.8$ percentage-point proportionality gap (3/8 = 37.5\,\% seats vs
50.3\,\% votes).

GeoSection creates two districts at exactly 49.4\,\% Democratic vote share ---
knife-edge seats that could flip with any modest swing.
These districts arise as a direct byproduct of urban peeling: when Milwaukee
is carved off as a single district, the suburban ring districts that border
it land near 50/50.

AreaSection creates zero near-toss-up seats.
The geographic balance forces Milwaukee's population to be distributed between
the two halves of the state rather than isolated in one, producing a cleaner
partisan sorting: safe Democratic districts (72.8\,\%, 68.5\,\%, 51.5\,\%) and
clearly Republican districts (42.7\,\%, 42.1\,\%, 38.4\,\%, 46.2\,\%).

\subsection{North Carolina: Head-to-Head Comparison}

\begin{table}[h]
\centering
\caption{Summary comparison: GeoSection vs AreaSection, North Carolina
(14 districts).}
\label{tab:nc-summary}
\begin{tabular}{lrr}
\toprule
Metric & GeoSection & AreaSection \\
\midrule
Natural ratio & 1:13 & 6:8 \\
Top-level edge cut & 97\,km & 279\,km \\
D seats / R seats & 5 / 9 & 5 / 9 \\
D vote share & 48.9\,\% & 48.9\,\% \\
Proportionality gap & $-7.1$\,pp & $-7.1$\,pp \\
\bottomrule
\end{tabular}
\end{table}

North Carolina reinforces the Wisconsin finding: same seat allocation, same
proportionality gap, different top-level geometry.
GeoSection peels Charlotte/Raleigh at 97\,km boundary; AreaSection forces a
6:8 east/west split at a 279\,km boundary --- $2.9\times$ longer.

The boundary length difference is the price of geographic balance.
Whether that price is worth paying is a legal and normative question, not an
algorithmic one.

\subsection{The Lorenz Curve as a State-Level Diagnostic}

Figure~\ref{fig:lorenz-wi} shows the population-area Lorenz curve for Wisconsin.
The curve reaches $(0.5, 0.932)$: the densest 50\,\% of Wisconsin's land holds
93.2\,\% of the population.
The vertical line at area fraction 0.5 intersects the curve far above the
diagonal, confirming severe population concentration.

For comparison, Iowa (uniform agriculture-based density) has a Lorenz curve
close to the diagonal, $p^* \approx 0.55$, and AreaSection selects a 2:2
equal split.

The Lorenz curve is a useful \emph{one-number summary} of whether a state is
likely to produce equal or asymmetric splits under AreaSection:
states with $p^* > 0.8$ (highly concentrated) retain asymmetric splits;
states with $p^* < 0.65$ (relatively uniform) produce equal or near-equal splits.
